\bta{2016}

\begin{enumerate}
	%\renewcommand{\labelenumi}{\arabic{enumi}.}
	% A(\Alph) a(\alph) I(\Roman) i(\roman) 1(\arabic)
	%设定全局标号series=example	%引用全局变量resume=example
	%[topsep=-0.3em,parsep=-0.3em,itemsep=-0.3em,partopsep=-0.3em]
	%可使用leftmargin调整列表环境左边的空白长度 [leftmargin=0em]
	\item
	选择题 (本题满分 50 分，每小题 5 分)
	\begin{enumerate}
		%\renewcommand{\labelenumi}{\arabic{enumi}.}
		% A(\Alph) a(\alph) I(\Roman) i(\roman) 1(\arabic)
		%设定全局标号series=example	%引用全局变量resume=example
		%[topsep=-0.3em,parsep=-0.3em,itemsep=-0.3em,partopsep=-0.3em]
		%可使用leftmargin调整列表环境左边的空白长度 [leftmargin=0em]
		\item
		对于函数 $f(x)$, 当 $x$ 是有理数时, $f(x)=\sin \pi x$ ；当 $x$ 为无理数时, $f(x)=0$, 则正确的结论是 \xzanswer{D}


\fourchoices
{$f(x)$ 在 $(-\infty,+\infty)$ 内无界}
{$f(x)$ 处处不连续}
{$f(x)$ 在 $(-\infty,+\infty)$ 上处处连续}
{$f(x)$ 在整数点上连续}

\item
设 $f^{\prime}\left(\sin ^{2} x\right)=\cos ^{2} x$, 则 $f(x)$ 等于 \xzanswer{C}


\fourchoices
{$x-\frac{1}{2} x^{2}$}
{$x+\frac{1}{2} x^{2}$}
{$x-\frac{1}{2} x^{2}-C$}
{$\frac{1}{2} x^{2}-x+C$}

\item
设 $y_{1}$ 和 $y_{2}$ 均为微分方程 $y^{\prime \prime}+2 y^{\prime}+y=2 x e^{-x}$ 的特解，则如下论断不正确的为 \xzanswer{D}


\fourchoices
{$\frac{1}{2} y_{1}+\frac{1}{2} y_{2}$ 可能形如 $\left(c_{1} x^{3}+c_{2} x^{2}\right) e^{-x}$, 其中 $c_{1}, c_{2}$ 为常数}
{$\frac{1}{2} y_{1}-\frac{1}{2} y_{2}$ 可能形如 $\left(c_{3} x+c_{4}\right) e^{-x}$, 其中 $c_{3}, c_{4}$ 为常数}
{$2 y_{1}+y_{2}$ 可能形如 $\left(c_{5} x^{3}+c_{6} x^{2}\right) e^{-x}$, 其中 $c_{5}, c_{6}$ 为常数}
{$2 y_{1}-y_{2}$ 可能形如 $\left(c_{7} x+c_{8}\right) e^{-x}$, 其中 $c_{7}, c_{8}$ 为常数}


\item
$ \int_{0}^{\pi} \frac{x \sin x}{1+\cos ^{2} x} d x= $ \xzanswer{C}


\fourchoices
{$\frac{\pi}{2}$}
{$\frac{\pi}{4}$}
{$\frac{\pi^{2}}{4}$}
{$\frac{\pi^{2}}{6}$}

\item
已知实向量 $a, b$ 满足 $(a+b) \perp(a-b),|a+b|=1,|a|=1$, 则 $|a \times b|=$ \xzanswer{A}


\fourchoices
{$\frac{\sqrt{3}}{2}$}
{$ 1 $}
{$\frac{\sqrt{2}}{2}$}
{$\frac{1}{6}$}

\item
设 $f(x)$ 为连续函数, $F(t)=\int_{1}^{t} d y \int_{y}^{t} f(x) d x$, 则 $F^{\prime}(2)=$ \xzanswer{B}


\fourchoices
{$2 f(0)$}
{$f(2)$}
{$-f(2)$}
{$ 0 $}

\item
正项级数 $\sum\limits_{n=1}^{+\infty} a_{n} $ 收敛是级数 $\sum\limits_{n=1}^{+\infty} a_{n}^{2} $ 收敛的 \xzanswer{B}


\fourchoices
{充要条件}
{充分条件，但非必要条件}
{必要条件，但非充分条件}
{既非充分条件，也非必要条件}

\item
$ \lim\limits _{x \rightarrow 0, y \rightarrow 0} \frac{x^{2} y}{\sqrt[3]{x^{8}+y^{12}}}= $ \xzanswer{D}


\fourchoices
{$ 0 $}
{$\frac{1}{\sqrt[3]{3}}$}
{$\frac{1}{\sqrt[3]{2}}$}
{不存在}

\item
求 $\lim\limits _{x \rightarrow 0}\left(\frac{\sin x}{x}\right)^{\frac{1}{x^{2}}}=$ \xzanswer{B}


\fourchoices
{$e^{\frac{1}{6}}$}
{$e^{-\frac{1}{6}}$}
{$e^{\frac{1}{2}}$}
{$e^{-\frac{1}{2}}$}

\item
若直线 $\left\{\begin{array}{l}x-y+z-1=0 \\ 2 x+y=0\end{array}\right.$ 与平面 $x+k y-z-5=0$ 平行，则 $k$ 的取值为 \xzanswer{C}


\fourchoices
{$ -1 $}
{$ 0 $}
{$ 2 $}
{$ 1 $}





	\end{enumerate}


\banswer{
\begin{enumerate}
	%\renewcommand{\labelenumi}{\arabic{enumi}.}
	% A(\Alph) a(\alph) I(\Roman) i(\roman) 1(\arabic)
	\item
	\textbf{D}。对一切 $x$ 都有 $|f(x)|\leqslant|\sin\pi x|\leqslant1$，故 $f$ 有界，A 错。
	设 $x_0$ 为整数，则 $x_0$ 是有理数且 $f\left(x_0\right)=\sin\pi x_0=0$；
	由 $|f(x)-f\left(x_0\right)|\leqslant|\sin\pi x|\to0\ (x\to x_0)$ 知 $f$ 在整数点连续，
	D 对而 B 错。若 $x_0$ 不是整数，则 $\sin\pi x_0\neq0$：
	$x_0$ 的任一邻域内既有有理点（$f$ 的值接近 $\sin\pi x_0\neq0$）
	又有无理点（$f=0$），故 $f$ 在非整数点必不连续，C 错。

	\item
	\textbf{C}。令 $t=\sin^2x$，则 $\cos^2x=1-t$，且 $t$ 取遍 $[0,1]$，故 $f'(t)=1-t$，
	\[
	f(t)=t-\frac{t^2}{2}+C .
	\]
	$C$ 是任意常数，$-C$ 同样任意，故答案形如 $x-\frac12x^2-C$，即 C。
	A 漏掉了任意常数，只是 $C=0$ 的一个特例；B、D 的系数或符号都不对。

	\item
	\textbf{D}。特征方程 $r^2+2r+1=0$ 有二重根 $r=-1$。令 $y=v(x)e^{-x}$，
	直接计算得 $y''+2y'+y=v''e^{-x}$，故原方程等价于 $v''=2x$，
	$v=\frac{x^3}{3}+\alpha x+\beta$。因此原方程的任一特解必形如
	\[
	y=\left(\frac{x^3}{3}+\alpha x+\beta\right)e^{-x} .
	\]
	线性组合的系数之和决定它满足哪个方程：

	A：系数和 $\frac12+\frac12=1$，仍是原方程的解，取 $\alpha=\beta=0$ 即
	$\frac13x^3e^{-x}$，可写成 $\left(c_1x^3+c_2x^2\right)e^{-x}$（$c_2=0$），论断成立；

	B：系数和 $0$，是齐次方程的解，即 $\left(c_3x+c_4\right)e^{-x}$，论断成立；

	C：系数和 $3$，满足 $y''+2y'+y=6xe^{-x}$，其解形如
	$\left(x^3+\alpha x+\beta\right)e^{-x}$，取 $\alpha=\beta=0$ 即可，论断成立；

	D：系数和 $1$，仍是原方程的解，$x^3$ 的系数恒为 $\frac13\neq0$，
	绝不可能退化成 $\left(c_7x+c_8\right)e^{-x}$，故 D 不正确。

	\item
	\textbf{C}。作代换 $x=\pi-u$ 可证：若 $g$ 只依赖 $\sin x$，则
	$\int_0^{\pi}xg(\sin x)dx=\frac{\pi}{2}\int_0^{\pi}g(\sin x)dx$。
	本题 $\frac{\sin x}{1+\cos^2x}=\frac{\sin x}{2-\sin^2x}$ 正是这种函数，故
	\[
	I=\frac{\pi}{2}\int_0^{\pi}\frac{\sin x}{1+\cos^2x}dx
	=\frac{\pi}{2}\Big[-\arctan(\cos x)\Big]_0^{\pi}=\frac{\pi^2}{4} .
	\]

	\item
	\textbf{A}。由 $(\bm a+\bm b)\cdot(\bm a-\bm b)=|\bm a|^2-|\bm b|^2=0$ 得 $|\bm b|=|\bm a|=1$；
	再由 $1=|\bm a+\bm b|^2=|\bm a|^2+2\bm a\cdot\bm b+|\bm b|^2=2+2\bm a\cdot\bm b$
	得 $\bm a\cdot\bm b=-\frac12$。故夹角满足 $\cos\theta=-\frac12$，$\sin\theta=\frac{\sqrt3}{2}$，
	$|\bm a\times\bm b|=|\bm a||\bm b|\sin\theta=\frac{\sqrt3}{2}$。

	\item
	\textbf{B}。交换积分次序：积分区域
	$\{1\leqslant y\leqslant t,\ y\leqslant x\leqslant t\}$
	即 $\{1\leqslant x\leqslant t,\ 1\leqslant y\leqslant x\}$，故
	\[
	F(t)=\int_1^{t}f(x)\left(\int_1^{x}dy\right)dx=\int_1^{t}(x-1)f(x)dx .
	\]
	于是 $F'(t)=(t-1)f(t)$，$F'(2)=1\cdot f(2)=f(2)$。

	\item
	\textbf{B}。充分性：$\sum a_n$ 收敛 $\Rightarrow a_n\to0$，故存在 $N$ 使 $n>N$ 时
	$0\leqslant a_n<1$，从而 $a_n^2<a_n$，由比较判别法 $\sum a_n^2$ 收敛。
	非必要：取 $a_n=\frac1n$，则 $\sum a_n^2=\sum\frac1{n^2}$ 收敛而 $\sum a_n$ 发散。

	\item
	\textbf{D}。沿 $y=x$（$x\to0^{+}$）：
	$\frac{x^3}{\sqrt[3]{x^8+x^{12}}}\sim\frac{x^3}{x^{8/3}}=x^{1/3}\to0$；
	沿 $y=x^{2/3}$（$x\to0^{+}$）：
	\[
	\frac{x^2\cdot x^{2/3}}{\sqrt[3]{x^8+x^8}}=\frac{x^{8/3}}{\sqrt[3]{2}\,x^{8/3}}=\frac{1}{\sqrt[3]{2}} .
	\]
	两条路径给出不同的值，故二重极限不存在。（C 只是沿某一条路径的值。）

	\item
	\textbf{B}。由 $\frac{\sin x}{x}=1-\frac{x^2}{6}+o\left(x^2\right)$ 及 $\ln(1+u)\sim u$，
	\[
	\frac{1}{x^2}\ln\frac{\sin x}{x}=\frac{1}{x^2}\left(-\frac{x^2}{6}+o\left(x^2\right)\right)\to-\frac16 ,
	\]
	故原极限为 $e^{-\frac16}$。

	\item
	\textbf{C}。直线的方向向量取两平面法向的叉积
	\[
	\bm s=(1,-1,1)\times(2,1,0)=(-1,\,2,\,3) ,
	\]
	平面法向 $\bm n=(1,k,-1)$。直线与平面平行要求 $\bm s\cdot\bm n=0$，即
	$-1+2k-3=0$，$k=2$。再验证直线不落在平面内：
	取线上一点 $(0,0,1)$，代入平面左端得 $-6\neq0$，故确为平行。
\end{enumerate}

\par\vspace{1mm}\noindent
\textbf{注：}{\kaishu
第 9 小题原题面印作 $\left(\frac{\sin}{x}\right)^{1/x^2}$，分子脱落了自变量，
已按通行版本补成 $\left(\frac{\sin x}{x}\right)^{1/x^2}$；
第 3 小题 D 选项原印作 $\left(c_7x+c_4\right)e^{-x}$，与其后“$c_7,c_8$ 为常数”不符，
已改为 $\left(c_7x+c_8\right)e^{-x}$。两处都不影响答案。
}\par
}



\item
(本题满分 10 分)
若 $\lim\limits _{n \rightarrow \infty} x_{n}=a, \lim\limits _{n \rightarrow \infty} y_{n}=b$, 计算下列极限
\[
\lim\limits _{n \rightarrow \infty} \frac{x_{1} y_{n}+x_{2} y_{n-1}+\cdots+x_{n} y_{1}}{n}
\]
并写出具体过程。
\banswer{
	所求极限为 $ab$。

	\textbf{引理（Cauchy 命题）.} 若 $\gamma_n\to0$，则
	$\frac{\gamma_1+\cdots+\gamma_n}{n}\to0$。
	事实上，任给 $\varepsilon>0$，取 $N$ 使 $k>N$ 时 $|\gamma_k|<\frac{\varepsilon}{2}$，则
	\[
	\left|\frac{1}{n}\sum_{k=1}^{n}\gamma_k\right|
	\leqslant\frac{1}{n}\sum_{k=1}^{N}\left|\gamma_k\right|+\frac{\varepsilon}{2} ,
	\]
	右端第一项当 $n$ 充分大时 $<\frac{\varepsilon}{2}$（分子是固定的数）。

	\textbf{主证.} 记 $x_n=a+\alpha_n$，$y_n=b+\beta_n$，则 $\alpha_n\to0$，$\beta_n\to0$。于是
	\begin{align*}
	\frac{1}{n}\sum_{k=1}^{n}x_ky_{n+1-k}
	=\ &ab+\frac{b}{n}\sum_{k=1}^{n}\alpha_k+\frac{a}{n}\sum_{k=1}^{n}\beta_k\\
	&+\frac{1}{n}\sum_{k=1}^{n}\alpha_k\beta_{n+1-k} .
	\end{align*}
	由引理，第二、三项均 $\to0$。对第四项：收敛数列有界，设 $|\alpha_k|\leqslant M$，则
	\[
	\left|\frac{1}{n}\sum_{k=1}^{n}\alpha_k\beta_{n+1-k}\right|
	\leqslant\frac{M}{n}\sum_{j=1}^{n}\left|\beta_j\right|\to0 ,
	\]
	再次用到引理（对 $\left|\beta_j\right|\to0$）。综上，原极限 $=ab$。
}


\item
(本题满分 10 分)
求满足下列初始值的常微分方程 $y=y(x)$：
\[
\left\{\begin{aligned}
	&y^{\prime \prime}=y^{\prime}(2 y+2) \\
	&y(0)=-1 \\
	&y^{\prime}(0)=1
\end{aligned}\right.
\]
的解。

\banswer{
	方程不显含 $x$，令 $p=y'$ 并把 $p$ 看成 $y$ 的函数，则
	$y''=\frac{dp}{dx}=p\frac{dp}{dy}$，原方程化为
	\[
	p\frac{dp}{dy}=p(2y+2) .
	\]

	\textbf{第一次积分.} 初值给出 $p(0)=1\neq0$，故在初值附近可约去 $p$：
	$\frac{dp}{dy}=2y+2$，积分得 $p=y^2+2y+C_1$。
	由 $y=-1$ 时 $p=1$ 得 $1=1-2+C_1$，$C_1=2$，于是
	\[
	\frac{dy}{dx}=y^2+2y+2=(y+1)^2+1 .
	\]

	\textbf{第二次积分.} 令 $u=y+1$，则 $\frac{du}{1+u^2}=dx$，积分得
	$\arctan u=x+C_2$。由 $x=0$ 时 $u=0$ 得 $C_2=0$，故 $u=\tan x$，即
	\[
	y=\tan x-1 ,\qquad x\in\left(-\frac{\pi}{2},\frac{\pi}{2}\right) .
	\]
	（验证：$y'=\sec^2x$，$y''=2\sec^2x\tan x$，而
	$y'(2y+2)=\sec^2x\cdot2\tan x$，相等。）

	\par\vspace{1mm}\noindent
	\textbf{注：}{\kaishu
	另有解 $p\equiv0$（即 $y\equiv$ 常数），但它不满足 $y'(0)=1$，故舍去。
	解的存在区间只能到 $x=\pm\frac{\pi}{2}$，在那里 $y$ 爆破。
	}\par
}


\item
(本题满分 10 分)
一空间物体由球面 $(x-1)^{2}+y^{2}+z^{2}=1$ 的内侧和锥面 $x=\sqrt{y^{2}+z^{2}}$ 朝向 $x$ 轴 正向的一侧所界定, 其在点 $(x, y, z)$ 的密度为 $\rho(x, y, z)=1-\left(y^{2}+z^{2}\right)$, 试求该物体的质量。

\banswer{
	物体关于 $x$ 轴旋转对称，故把 $x$ 轴取作极轴用球坐标：
	\[
	x=r\cos\varphi,\quad y=r\sin\varphi\cos\theta,\quad z=r\sin\varphi\sin\theta ,
	\]
	$dV=r^2\sin\varphi\,dr\,d\varphi\,d\theta$，且 $y^2+z^2=r^2\sin^2\varphi$。

	\textbf{定区域.} 球面 $(x-1)^2+y^2+z^2=1$ 即 $x^2+y^2+z^2=2x$，
	化为 $r=2\cos\varphi$；锥面 $x=\sqrt{y^2+z^2}$ 即 $\cos\varphi=\sin\varphi$，
	即 $\varphi=\frac{\pi}{4}$。锥内（朝 $x$ 轴正向的一侧）对应 $0\leqslant\varphi\leqslant\frac{\pi}{4}$，故
	\[
	\Omega:\ 0\leqslant\theta<2\pi,\ \ 0\leqslant\varphi\leqslant\frac{\pi}{4},\ \ 0\leqslant r\leqslant2\cos\varphi .
	\]

	\textbf{积分.} 密度 $\rho=1-r^2\sin^2\varphi$，先对 $r$ 积分：
	\[
	M=2\pi\int_0^{\pi/4}\!\!\sin\varphi
	\left[\frac{r^3}{3}-\frac{r^5\sin^2\varphi}{5}\right]_0^{2\cos\varphi}\!\!d\varphi .
	\]
	令 $c=\cos\varphi$（$dc=-\sin\varphi\,d\varphi$，$\sin^2\varphi=1-c^2$）：
	\[
	M=2\pi\int_{\frac{\sqrt2}{2}}^{1}\left[\frac{8c^3}{3}-\frac{32c^5\left(1-c^2\right)}{5}\right]dc .
	\]
	原函数为 $\frac{2c^4}{3}-\frac{16c^6}{15}+\frac{4c^8}{5}$，在 $c=1$ 处的值是 $\frac25$，
	在 $c=\frac{\sqrt2}{2}$ 处是 $\frac16-\frac{2}{15}+\frac{1}{20}=\frac{1}{12}$，故
	\[
	M=2\pi\left(\frac25-\frac{1}{12}\right)=2\pi\cdot\frac{19}{60}=\frac{19}{30}\pi .
	\]

	\par\vspace{1mm}\noindent
	\textbf{注：}{\kaishu
	锥面把球分成两块，本题取的是含球心、朝 $x$ 轴正向张开的那一块
	（即 $\varphi\leqslant\frac{\pi}{4}$ 的部分），不要误取补集。
	蒙特卡洛数值积分给出 $M\approx1.9890$，而 $\frac{19}{30}\pi\approx1.98968$，吻合。
	}\par
}


\item
(本题满分 10 分)
将函数 $f(x)=(\pi-|x|)^{2}(-\pi \leqslant x \leqslant \pi)$ 展开成傅里叶级数，并求 $\sum\limits_{n=1}^{+\infty} \frac{1}{n^{2}}$ 的和。

\banswer{
	$f$ 是偶函数，故 $b_n=0$，只需算余弦系数（在 $[0,\pi]$ 上 $f(x)=(\pi-x)^2$）。

	\textbf{定 $a_0$.} $a_0=\frac{2}{\pi}\int_0^{\pi}(\pi-x)^2dx=\frac{2}{\pi}\cdot\frac{\pi^3}{3}=\frac{2\pi^2}{3}$，
	常数项 $\frac{a_0}{2}=\frac{\pi^2}{3}$。

	\textbf{定 $a_n\ (n\geqslant1)$.} 令 $u=\pi-x$，注意 $\cos\left[n(\pi-u)\right]=(-1)^n\cos nu$：
	\[
	a_n=\frac{2(-1)^n}{\pi}\int_0^{\pi}u^2\cos nu\,du .
	\]
	两次分部积分得 $\int_0^{\pi}u^2\cos nu\,du=\frac{2\pi(-1)^n}{n^2}$，故 $a_n=\frac{4}{n^2}$。
	于是
	\[
	f(x)=\frac{\pi^2}{3}+\sum_{n=1}^{\infty}\frac{4}{n^2}\cos nx ,
	\qquad -\pi\leqslant x\leqslant\pi .
	\]
	（$f$ 连续且 $f(-\pi)=f(\pi)=0$，故周期延拓后处处连续，级数处处收敛于 $f$。）

	\textbf{求和.} 取 $x=0$，$f(0)=\pi^2$：
	\[
	\pi^2=\frac{\pi^2}{3}+4\sum_{n=1}^{\infty}\frac{1}{n^2}
	\ \Longrightarrow\
	\sum_{n=1}^{\infty}\frac{1}{n^2}=\frac{1}{4}\cdot\frac{2\pi^2}{3}=\frac{\pi^2}{6} .
	\]

	\par\vspace{1mm}\noindent
	\textbf{注：}{\kaishu
	若改取 $x=\pi$，得 $0=\frac{\pi^2}{3}+4\sum\frac{(-1)^n}{n^2}$，
	即 $\sum\frac{(-1)^{n-1}}{n^2}=\frac{\pi^2}{12}$，两式相容。
	}\par
}


\item
(本题满分 10 分)
设函数 $f(x)=\int_{0}^{x} g(t) d t$, 求 $\int_{0}^{1} e^{-x^{2}} f(x) d x$ 的值, 其中 $g(t)=5 t^{4}+3 t^{2}+1$。

\banswer{
	先把 $f$ 算出来：
	\[
	f(x)=\int_0^{x}\left(5t^4+3t^2+1\right)dt=x^5+x^3+x=x\left(x^4+x^2+1\right) .
	\]
	被积函数中 $x$ 的奇次幂提示换元 $u=x^2$（$du=2x\,dx$）：
	\[
	\int_0^1e^{-x^2}f(x)dx=\frac12\int_0^1e^{-u}\left(u^2+u+1\right)du .
	\]
	对多项式 $P$ 有 $\int e^{-u}P\,du=-e^{-u}\left(P+P'+P''+\cdots\right)$。
	取 $P=u^2+u+1$，则 $P+P'+P''=u^2+3u+4$，故
	\[
	\int_0^1e^{-u}\left(u^2+u+1\right)du
	=\Big[-e^{-u}\left(u^2+3u+4\right)\Big]_0^1=4-\frac{8}{e} .
	\]
	所以
	\[
	\int_0^1e^{-x^2}f(x)dx=\frac12\left(4-\frac8e\right)=2-\frac{4}{e}\approx0.5285 .
	\]
}


\newpage
\item
(本题满分 10 分)
设在上半平面 $D=\{(x, y) \mid y>0\}$ 内, 函数 $f(x, y)$ 具有连续偏导数, 且对任意的 $t>0$ 都有 $f(t x, t y)=t^{-2} f(x, y)$ 。 证明：对 $D$ 内的任意分段光滑的有向简单闭曲线 $L$, 都有
\[
\oint_{L} y f(x, y) d x-x f(x, y) d y=0
\]

\banswer{
	\textbf{第一步：欧拉恒等式.} 把 $f(tx,ty)=t^{-2}f(x,y)$ 两端对 $t$ 求导：
	\[
	xf_x(tx,ty)+yf_y(tx,ty)=-2t^{-3}f(x,y) ,
	\]
	令 $t=1$ 即得 $-2$ 次齐次函数的欧拉恒等式
	\[
	xf_x(x,y)+yf_y(x,y)=-2f(x,y) .
	\]

	\textbf{第二步：格林公式.} 记 $P=yf$，$Q=-xf$，二者在 $D$ 上有连续偏导数，且
	\begin{align*}
	\frac{\partial Q}{\partial x}-\frac{\partial P}{\partial y}
	&=\left(-f-xf_x\right)-\left(f+yf_y\right)\\
	&=-2f-\left(xf_x+yf_y\right)=-2f+2f=0 .
	\end{align*}
	设 $L$ 所围的区域为 $D_L$。上半平面是单连通的，故 $D_L\subset D$，
	在 $D_L$ 上可用格林公式：
	\[
	\oint_{L}yf\,dx-xf\,dy=\pm\iint_{D_L}\left(Q_x-P_y\right)dxdy=0
	\]
	（正负号取决于 $L$ 的定向，不影响结论为零）。

	\par\vspace{1mm}\noindent
	\textbf{注：}{\kaishu
	这正是力学中位力定理所用的技巧：$k$ 次齐次函数满足
	$xf_x+yf_y=kf$，本题 $k=-2$ 恰好抵消了格林公式里多出来的 $-2f$。
	}\par
}


\item
(本题满分 10 分)
计算曲面积分
\[
I=\iint_{\Sigma} x z d y d z+2 y z d z d x+3 x y d x d y
\]
其中, $\Sigma$ 为曲面 $z=1-x^{2}-\frac{1}{4} y^{2}$ $(0 \leqslant z \leqslant 1)$ 的上侧。

\banswer{
	\textbf{解法一（高斯公式）.} 补上底面
	$\Sigma_1:\ z=0,\ x^2+\frac{y^2}{4}\leqslant1$ 并取下侧，
	则 $\Sigma+\Sigma_1$ 是所围立体 $\Omega$ 的外侧闭曲面，而
	\[
	\operatorname{div}\bm F=z+2z+0=3z .
	\]
	$\Omega$ 在高度 $z$ 处的截面是椭圆 $x^2+\frac{y^2}{4}\leqslant1-z$，
	半轴为 $\sqrt{1-z}$ 与 $2\sqrt{1-z}$，面积 $2\pi(1-z)$，故
	\[
	\iiint_{\Omega}3z\,dV=\int_0^1 3z\cdot2\pi(1-z)\,dz
	=6\pi\left(\frac12-\frac13\right)=\pi .
	\]
	在 $\Sigma_1$ 上 $dydz=dzdx=0$，只剩 $\iint3xy\,dxdy$，
	而积分域关于 $y$ 轴对称、$3xy$ 关于 $x$ 为奇函数，故该项为 $0$。于是
	\[
	I=\pi-0=\pi .
	\]

	\textbf{解法二（直接投影）.} 取上侧时
	$d\bm S=\left(-z_x,-z_y,1\right)dxdy=\left(2x,\frac{y}{2},1\right)dxdy$，故
	\[
	I=\iint_{D}\left[\left(2x^2+y^2\right)z+3xy\right]dxdy ,
	\]
	其中 $D:\ x^2+\frac{y^2}{4}\leqslant1$，$z=1-x^2-\frac{y^2}{4}$。
	$3xy$ 一项由对称性为零。作广义极坐标
	$x=\rho\cos\theta$，$y=2\rho\sin\theta$（$dxdy=2\rho\,d\rho\,d\theta$），
	则 $z=1-\rho^2$，$2x^2+y^2=\rho^2\left(3-\cos2\theta\right)$，于是
	\[
	I=\int_0^{2\pi}\!\!\left(3-\cos2\theta\right)d\theta
	\int_0^1\!\!\rho^2\left(1-\rho^2\right)2\rho\,d\rho
	=6\pi\times\frac16=\pi ,
	\]
	与解法一一致（数值积分亦给出 $I\approx3.14159$）。
}


\item
(本题满分 10 分)
函数 $f(x)$ 在 $(0,\infty)$ 上可导, 且 $\lim\limits _{x \rightarrow \infty} f^{\prime}(x)=0$, 证明:
\[
\lim\limits _{x \rightarrow \infty} \frac{f(x)}{x}=0
\]

\banswer{
	\textbf{解法一（拉格朗日中值定理）.} 任给 $\varepsilon>0$，由 $f'(x)\to0$，
	存在 $X>0$ 使 $x>X$ 时 $\left|f'(x)\right|<\frac{\varepsilon}{2}$。
	对 $x>X$，在 $[X,x]$ 上用中值定理：存在 $\xi\in(X,x)$ 使
	\[
	f(x)-f(X)=f'(\xi)(x-X) ,
	\]
	故 $|f(x)|\leqslant|f(X)|+\frac{\varepsilon}{2}(x-X)\leqslant|f(X)|+\frac{\varepsilon}{2}x$，从而
	\[
	\left|\frac{f(x)}{x}\right|\leqslant\frac{|f(X)|}{x}+\frac{\varepsilon}{2} .
	\]
	$|f(X)|$ 是与 $x$ 无关的定数，故存在 $X_1>X$ 使 $x>X_1$ 时
	$\frac{|f(X)|}{x}<\frac{\varepsilon}{2}$，于是 $\left|\frac{f(x)}{x}\right|<\varepsilon$。
	由 $\varepsilon$ 的任意性即得结论。

	\textbf{解法二（洛必达法则的 $\frac{*}{\infty}$ 型）.} 分母 $x\to+\infty$，
	且 $\lim\limits_{x\to\infty}\frac{f'(x)}{1}=0$ 存在，
	故可直接用洛必达法则（此型不要求分子趋于无穷）：
	\[
	\lim_{x\to\infty}\frac{f(x)}{x}=\lim_{x\to\infty}f'(x)=0 .
	\]

	\par\vspace{1mm}\noindent
	\textbf{注：}{\kaishu
	不能写成 $f(x)=\int_0^{x}f'(t)dt+f(0)$：$f$ 只在 $(0,\infty)$ 上有定义，
	$f(0)$ 未必存在，$f'$ 也未必可积。解法一只用到中值定理，条件最省。
	}\par
}


\item
(本题满分 10 分)
已知函数 $f(x)$ 在 $(-\infty,+\infty)$ 上连续, 且 $F(x)=\frac{1}{2 \delta} \int_{-\delta}^{\delta} f(x+t) d t$, 其中 $\delta>0$。 证明: $F^{\prime}(x)$ 存在且连续, 其中 $x \in(-\infty,+\infty)$。


\banswer{
	\textbf{化为变限积分.} 在积分中作代换 $u=x+t$（$x$ 固定，$du=dt$），积分限
	$t=\mp\delta$ 对应 $u=x\mp\delta$，故
	\[
	F(x)=\frac{1}{2\delta}\int_{x-\delta}^{x+\delta}f(u)\,du .
	\]

	\textbf{求导.} 因 $f$ 在 $\mathbb{R}$ 上连续，令 $G(x)=\int_0^{x}f(u)du$，
	由微积分基本定理 $G$ 可导且 $G'=f$。于是
	\[
	F(x)=\frac{1}{2\delta}\left[G(x+\delta)-G(x-\delta)\right] ,
	\]
	它是两个可导函数之差，故 $F$ 处处可导，且
	\[
	F'(x)=\frac{f(x+\delta)-f(x-\delta)}{2\delta} .
	\]

	\textbf{连续性.} 上式右端是连续函数 $f$ 在两个连续变动点处取值的线性组合，
	故 $F'$ 在 $(-\infty,+\infty)$ 上连续。

	\par\vspace{1mm}\noindent
	\textbf{注：}{\kaishu
	注意 $\delta$ 是常数而非变量，所以求导只作用在积分限的 $x$ 上。
	$F$ 是 $f$ 在长度 $2\delta$ 的窗口上的滑动平均，此结论说明“磨光”确实
	把连续函数提升为连续可导函数。
	}\par
}


\item
(本题满分 10 分)
设函数 $f(x)$ 在 $[a, b]$ 上二阶可导, $f^{\prime}(a)=f^{\prime}(b)=0$ 。证明：存在一点 $\xi \in(a, b)$, 使 得
\[
\left|f^{\prime \prime}(\xi)\right| \geq \frac{4}{(b-a)^{2}}|f(b)-f(a)|
\]

	\banswer{
		记 $c=\frac{a+b}{2}$。分别在 $a$、$b$ 处把 $f$ 展成带拉格朗日余项的一阶泰勒公式
		并在 $x=c$ 处取值，注意 $f'(a)=f'(b)=0$、$c-a=-(c-b)=\frac{b-a}{2}$：
		\begin{align*}
		f(c)&=f(a)+\frac{f''\left(\xi_1\right)}{2}\cdot\frac{(b-a)^2}{4},
		&&\xi_1\in(a,c),\\
		f(c)&=f(b)+\frac{f''\left(\xi_2\right)}{2}\cdot\frac{(b-a)^2}{4},
		&&\xi_2\in(c,b).
		\end{align*}
		两式相减得
		\[
		f(b)-f(a)=\frac{(b-a)^2}{8}\left[f''\left(\xi_1\right)-f''\left(\xi_2\right)\right] .
		\]
		取 $\xi$ 为 $\xi_1,\xi_2$ 中使 $\left|f''\right|$ 较大者，则
		\[
		|f(b)-f(a)|\leqslant\frac{(b-a)^2}{8}\cdot2\left|f''(\xi)\right|
		=\frac{(b-a)^2}{4}\left|f''(\xi)\right| ,
		\]
		即 $\left|f''(\xi)\right|\geqslant\frac{4}{(b-a)^2}|f(b)-f(a)|$，且 $\xi\in(a,b)$。
	}



\end{enumerate}


